\chapter{Relación ingresos-gastos y viabilidad económica}
\label{anexo:relacion_ingresos_gastos}

Este anexo recoge la relación entre los recursos económicos disponibles, los gastos asociados al prototipo TierraViva y una estimación del coste total del proyecto. El objetivo es analizar la viabilidad económica de un prototipo académico funcional, valorando tanto el gasto material necesario para reproducirlo como el esfuerzo técnico asociado a su desarrollo y su posible evolución hacia una solución más completa de riego inteligente.

Los costes materiales detallados de cada componente se recogen en el Anexo \ref{anexo:materiales}. En este apartado se ofrece una visión agregada, diferenciando entre gasto directo del prototipo, recursos disponibles, coste estimado de desarrollo y viabilidad de reproducción.

\section{Criterio económico utilizado}

Para analizar el proyecto de forma realista se distinguen tres conceptos:

\begin{itemize}
    \item \textbf{Gasto directo}: dinero destinado a la compra de componentes, materiales y servicios necesarios para construir el prototipo.
    \item \textbf{Recursos no monetizados}: recursos utilizados en el proyecto sin traducirse en un pago directo, como tutorización académica, software libre o acceso a un entorno de pruebas.
    \item \textbf{Coste valorado}: estimación económica del tiempo dedicado al análisis, diseño, montaje, programación, pruebas y documentación.
\end{itemize}

Esta separación permite distinguir el coste material de construcción del prototipo y el valor económico completo del trabajo realizado.

\section{Resumen económico del prototipo}

A partir del anexo de materiales, el coste directo estimado del prototipo asciende a \textbf{420.50 \euro{}}. Este importe incluye dos estaciones de campo, el nodo de supervisión con Raspberry Pi, comunicaciones, alimentación, actuación y material auxiliar de montaje.

\begin{table}[H]
\centering
\renewcommand{\arraystretch}{1.25}
\small
\begin{tabular}{|p{0.45\textwidth}|p{0.25\textwidth}|p{0.20\textwidth}|}
\hline
\rowcolor{URJCRed}
\TVHead{Concepto} & \TVHead{Tipo} & \TVHead{Importe} \\
\hline
Coste material del prototipo & Gasto directo & 420.50 \euro{} \\
\hline
Reserva para imprevistos & 10\% del coste material & 42.05 \euro{} \\
\hline
\rowcolor{gray!15}
\textbf{Presupuesto recomendado de reproducción} & \textbf{Material + imprevistos} & \textbf{462.55 \euro{}} \\
\hline
\end{tabular}
\caption{Resumen económico directo del prototipo TierraViva.}
\label{tab:resumen_economico_directo}
\end{table}

La reserva de imprevistos representa un margen razonable para cubrir sustitución de componentes, cables adicionales, conectores, material de aislamiento o pequeñas compras surgidas durante el montaje.

\section{Distribución de gastos por bloques}

Para facilitar la lectura económica del prototipo, los componentes se agrupan por bloques funcionales. Esta agrupación permite identificar qué partes del sistema concentran mayor inversión.

\begin{table}[H]
\centering
\renewcommand{\arraystretch}{1.25}
\small
\begin{tabular}{|p{0.32\textwidth}|p{0.43\textwidth}|p{0.15\textwidth}|}
\hline
\rowcolor{URJCRed}
\TVHead{Bloque} & \TVHead{Elementos incluidos} & \TVHead{Coste} \\
\hline
Sensorización & Sensores de humedad de suelo, BME280 y BH1750. & 16.40 \euro{} \\
\hline
Procesamiento y supervisión & Arduino UNO y Raspberry Pi 5. & 233.62 \euro{} \\
\hline
Comunicaciones & Módulos LTE A7670E-LASA y tarjetas SIM. & 38.00 \euro{} \\
\hline
Alimentación y actuación & Baterías, relé, LM2596 y bombas de agua. & 114.48 \euro{} \\
\hline
Montaje y material auxiliar & Protoboard, cableado, resistencias, condensadores y consumibles. & 18.00 \euro{} \\
\hline
\rowcolor{gray!15}
\textbf{Total} & \textbf{Coste material estimado del prototipo.} & \textbf{420.50 \euro{}} \\
\hline
\end{tabular}
\caption{Distribución de gastos del prototipo por bloques funcionales.}
\label{tab:gastos_por_bloques}
\end{table}

Como se observa en la Tabla \ref{tab:gastos_por_bloques}, el mayor peso económico corresponde al bloque de procesamiento y supervisión, principalmente por la Raspberry Pi 5. Este coste debe interpretarse correctamente: la Raspberry Pi actúa como nodo central del sistema y puede reutilizarse para varias estaciones siempre que la carga de datos y consulta se mantenga dentro de un rango asumible.

\section{Coste aproximado por estación}

El prototipo está planteado con dos estaciones remotas y un nodo central de supervisión. Si se excluye la Raspberry Pi 5, que funciona como recurso compartido, el coste aproximado asociado a las dos estaciones remotas es:

El coste asociado al conjunto de estaciones remotas se obtiene restando al coste total del prototipo el importe correspondiente al nodo central de supervisión:

\begin{center}
\small
\begin{tabular}{lr}
Coste total del prototipo & 420.50 \euro{} \\
Coste del nodo central Raspberry Pi & -204.20 \euro{} \\
\hline
Coste conjunto de estaciones remotas & 216.30 \euro{} \\
\end{tabular}
\end{center}

Al estar el prototipo formado por dos estaciones remotas, el coste orientativo por estación se estima dividiendo ese importe entre dos:

\begin{center}
\small
\begin{tabular}{lr}
Coste conjunto de estaciones remotas & 216.30 \euro{} \\
Número de estaciones & 2 \\
\hline
Coste orientativo por estación & 108.15 \euro{} \\
\end{tabular}
\end{center}

\begin{table}[H]
\centering
\renewcommand{\arraystretch}{1.25}
\small
\begin{tabular}{|p{0.48\textwidth}|p{0.25\textwidth}|p{0.17\textwidth}|}
\hline
\rowcolor{URJCRed}
\TVHead{Concepto} & \TVHead{Cálculo} & \TVHead{Importe} \\
\hline
Coste total del prototipo & Anexo de materiales & 420.50 \euro{} \\
\hline
Coste del nodo central Raspberry Pi & Raspberry Pi 5 & 204.20 \euro{} \\
\hline
Coste conjunto de estaciones remotas & 420.50 - 204.20 & 216.30 \euro{} \\
\hline
\rowcolor{gray!15}
\textbf{Coste orientativo por estación} & \textbf{216.30 / 2} & \textbf{108.15 \euro{}} \\
\hline
\end{tabular}
\caption{Estimación del coste orientativo por estación remota.}
\label{tab:coste_por_estacion}
\end{table}

Esta estimación es útil para analizar una posible ampliación futura del sistema. Aunque el coste exacto de una nueva estación dependería de precios, disponibilidad y componentes equivalentes, el cálculo muestra que el nodo central representa una inversión inicial compartida, mientras que cada estación adicional tendría un coste incremental menor.

\section{Ingresos y recursos disponibles}

Durante el desarrollo del TFG, el proyecto se ha financiado mediante recursos propios y apoyos académicos o experimentales. Por ello, en este anexo el término ingresos se interpreta como recursos económicos o no económicos que han permitido realizar el proyecto.

\begin{table}[H]
\centering
\renewcommand{\arraystretch}{1.25}
\small
\begin{tabular}{|p{0.26\textwidth}|p{0.43\textwidth}|p{0.19\textwidth}|}
\hline
\rowcolor{URJCRed}
\TVHead{Recurso} & \TVHead{Descripción} & \TVHead{Valor económico} \\
\hline
Aportación propia & Compra de componentes y materiales necesarios para construir el prototipo. & 420.50 \euro{} \\
\hline
Tutorización académica & Seguimiento técnico, orientación de ingeniería y contextualización experimental del proyecto. & No monetizado \\
\hline
Agrolab de Móstoles & Posibilidad de disponer de una parcela para realizar pruebas en un entorno más cercano al uso real. & No monetizado \\
\hline
Software libre, abierto o de uso gratuito en el alcance del prototipo & Arduino IDE, Python, FastAPI, SQLite, GitHub y tecnologías web. & 0 \euro{} \\
\hline
ThingSpeak & Plataforma IoT utilizada para recepción y consulta de telemetría durante el prototipo. & 0 \euro{} \\
\hline
\end{tabular}
\caption{Recursos económicos y no económicos disponibles para el desarrollo de TierraViva.}
\label{tab:ingresos_recursos_tierraviva}
\end{table}

La financiación mediante aportación propia y recursos académicos es coherente con la naturaleza del TFG. TierraViva se plantea como un prototipo funcional y documentado, orientado a validar una arquitectura técnica de riego inteligente y a servir como base para futuras evoluciones.

\section{Relación ingresos-gastos}

La relación económica directa del proyecto se resume en la Tabla \ref{tab:relacion_ingresos_gastos_final}. En ella se diferencia entre el balance real de caja y el presupuesto recomendado si otra persona quisiera reproducir el prototipo con un pequeño margen de seguridad.

\begin{table}[H]
\centering
\renewcommand{\arraystretch}{1.25}
\small
\begin{tabular}{|p{0.48\textwidth}|p{0.25\textwidth}|p{0.17\textwidth}|}
\hline
\rowcolor{URJCRed}
\TVHead{Concepto} & \TVHead{Tipo} & \TVHead{Importe} \\
\hline
Aportación propia para materiales & Recurso económico & 420.50 \euro{} \\
\hline
Coste material del prototipo & Gasto directo & -420.50 \euro{} \\
\hline
\rowcolor{gray!15}
\textbf{Balance económico directo} & \textbf{Ingresos - gastos} & \textbf{0.00 \euro{}} \\
\hline
Reserva recomendada para imprevistos & Margen de seguridad & -42.05 \euro{} \\
\hline
\rowcolor{gray!15}
\textbf{Presupuesto recomendado de reproducción} & \textbf{Material + 10\%} & \textbf{462.55 \euro{}} \\
\hline
Ingresos comerciales durante el TFG & No aplica & 0.00 \euro{} \\
\hline
\end{tabular}
\caption{Relación directa entre recursos económicos y gastos del prototipo.}
\label{tab:relacion_ingresos_gastos_final}
\end{table}

Desde el punto de vista económico directo, el proyecto queda equilibrado mediante aportación propia. Desde el punto de vista de reproducción, sería recomendable disponer de un presupuesto algo superior al coste material exacto para cubrir imprevistos habituales en prototipos electrónicos.

\section{Coste de desarrollo valorado}

Además del coste material, TierraViva requiere una dedicación significativa en análisis, diseño, programación, integración, pruebas y documentación. Estas horas representan el valor técnico del trabajo realizado dentro del marco académico del TFG.

Para realizar una estimación orientativa se toma una dedicación acumulada de \textbf{720 horas} entre finales de agosto de 2025, momento en el que comenzó la definición inicial del proyecto, y julio de 2026, fecha prevista para la entrega y defensa del TFG, teniendo en cuenta que los primeros meses estuvieron más centrados en planificación, revisión bibliográfica y definición de alcance, mientras que la mayor carga técnica y documental se concentró entre marzo y junio de 2026.

La estimación se plantea como una valoración razonable del esfuerzo técnico y documental asociado al desarrollo del prototipo. Incluye análisis, diseño, montaje, programación, integración, pruebas, validación en Agrolab Móstoles, elaboración de evidencias, redacción de capítulos y preparación de anexos.

Como coste interno de referencia se emplea un valor conservador de \textbf{15 \euro{}/h}. Esta cifra permite valorar de forma sencilla el esfuerzo técnico y documental asociado al proyecto.

\begin{table}[H]
\centering
\renewcommand{\arraystretch}{1.25}
\small
\begin{tabular}{|p{0.43\textwidth}|c|c|c|}
\hline
\rowcolor{URJCRed}
\TVHead{Actividad} & \TVHead{Horas} & \TVHead{Coste/h} & \TVHead{Coste} \\
\hline
Definición inicial de la idea y alcance & 40 h & 15 \euro{}/h & 600.00 \euro{} \\
\hline
Revisión de alternativas tecnológicas & 55 h & 15 \euro{}/h & 825.00 \euro{} \\
\hline
Reajuste de arquitectura y toma de decisiones técnicas & 50 h & 15 \euro{}/h & 750.00 \euro{} \\
\hline
Selección, compra y preparación de componentes & 35 h & 15 \euro{}/h & 525.00 \euro{} \\
\hline
Montaje, cableado e integración hardware & 70 h & 15 \euro{}/h & 1050.00 \euro{} \\
\hline
Pruebas del módulo LTE A7670E-LASA y comunicaciones & 55 h & 15 \euro{}/h & 825.00 \euro{} \\
\hline
Desarrollo del \textit{firmware} de estación & 80 h & 15 \euro{}/h & 1200.00 \euro{} \\
\hline
Integración con ThingSpeak & 35 h & 15 \euro{}/h & 525.00 \euro{} \\
\hline
Desarrollo en Raspberry Pi, API, SQLite y \textit{dashboard} & 85 h & 15 \euro{}/h & 1275.00 \euro{} \\
\hline
Pruebas, depuración y validación del sistema & 70 h & 15 \euro{}/h & 1050.00 \euro{} \\
\hline
Redacción de memoria, anexos y revisión final & 145 h & 15 \euro{}/h & 2175.00 \euro{} \\
\hline
\rowcolor{gray!15}
\textbf{Total desarrollo valorado} & \textbf{720 h} &  & \textbf{10800.00 \euro{}} \\
\hline
\end{tabular}
\caption{Estimación del coste de desarrollo valorado.}
\label{tab:coste_desarrollo_valorado}
\end{table}

\section{Coste total valorado del proyecto}

Si se considera únicamente el gasto material, TierraViva es un prototipo de bajo coste relativo. Sin embargo, si se valora el tiempo de desarrollo, integración, pruebas y documentación, el coste total del proyecto aumenta de forma significativa.

\begin{table}[H]
\centering
\renewcommand{\arraystretch}{1.25}
\small
\begin{tabular}{|p{0.52\textwidth}|p{0.25\textwidth}|}
\hline
\rowcolor{URJCRed}
\TVHead{Concepto} & \TVHead{Importe} \\
\hline
Coste material del prototipo & 420.50 \euro{} \\
\hline
Reserva de imprevistos & 42.05 \euro{} \\
\hline
Coste de desarrollo valorado & 10800.00 \euro{} \\
\hline
\rowcolor{gray!15}
\textbf{Coste total valorado con margen de reproducción} & \textbf{11262.55 \euro{}} \\\hline
\end{tabular}
\caption{Coste total valorado del proyecto TierraViva.}
\label{tab:coste_total_valorado}
\end{table}

Esta tabla permite distinguir claramente entre el coste material de reproducir el prototipo y el valor económico del trabajo técnico necesario para desarrollarlo. La diferencia entre ambos importes refleja que el valor de TierraViva procede de la combinación de componentes, diseño, integración, programación, pruebas y documentación.

\section{Escenarios de viabilidad}

Para analizar la viabilidad económica de TierraViva se plantean tres escenarios: prototipo académico, reproducción técnica y evolución futura.

\begin{table}[H]
\centering
\renewcommand{\arraystretch}{1.25}
\small
\begin{tabular}{|p{0.25\textwidth}|p{0.38\textwidth}|p{0.27\textwidth}|}
\hline
\rowcolor{URJCRed}
\TVHead{Escenario} & \TVHead{Descripción} & \TVHead{Valoración económica} \\
\hline
Prototipo académico & Sistema desarrollado para validar arquitectura, comunicaciones, sensorización, almacenamiento y actuación. & Viable con una inversión material aproximada de 420.50 \euro{}. \\
\hline
Reproducción técnica & Montaje equivalente por otra persona, considerando margen de seguridad para imprevistos. & Recomendable disponer de unos 462.55 \euro{}. \\
\hline
Ampliación con más estaciones & Incorporación de nuevas estaciones manteniendo el nodo central de supervisión. & Coste incremental orientativo de 108.15 \euro{} por estación, sujeto a variaciones de precios y disponibilidad de componentes. \\
\hline
Evolución precomercial & Mejora de robustez, cajas estancas, protección ambiental, seguridad, instalación y pruebas prolongadas. & Requeriría inversión adicional y una nueva estimación económica. \\
\hline
\end{tabular}
\caption{Escenarios de viabilidad económica del proyecto.}
\label{tab:escenarios_viabilidad}
\end{table}

\section{Conclusión}

El análisis económico muestra que TierraViva es viable como prototipo académico con un coste material contenido. La inversión directa se concentra principalmente en el nodo de supervisión y en los elementos de alimentación y actuación, mientras que los sensores y módulos de estación mantienen un coste relativamente bajo.

La principal diferencia aparece al valorar el tiempo de desarrollo. El valor principal del proyecto está en la integración de componentes, la programación del \textit{firmware}, la comunicación LTE, la conexión con ThingSpeak, el \textit{backend} en Raspberry Pi, la API, el \textit{dashboard}, las pruebas, la validación en Agrolab de Móstoles y la documentación técnica.

Por tanto, TierraViva puede considerarse un prototipo económicamente razonable para un Trabajo de Fin de Grado y, al mismo tiempo, una muy buena base técnica ampliable para futuras versiones próximas a una implantación real.